\section{Background}
\label{sec:background}

\paragraph{FAIR and its descendants.}
The FAIR principles~\cite{wilkinson2016fair} organise into fifteen
sub-principles (F1--F4, A1--A2, I1--I3, R1.1--R1.3) and were drafted
as a guideline set rather than a specification. FAIR4RS re-reads them
for research software, picking up versioning, dependencies, and
build reproducibility~\cite{chuehong2022fair4rs}.
FAIR4ML re-reads them for machine-learning models, where ``reusable''
gathers training-data documentation and intended-use statements
derived from datasheets~\cite{gebru2021datasheets} and model
cards~\cite{mitchell2019modelcards,ravi2024fair4ml}. Both extensions
are additive, not revisionist.
F(AI)\textsuperscript{2}R is the third such extension. It differs in
that the artefacts it covers are processual: a prompt is not the
output of research, but a versioned input that determines how the
output came about.

\paragraph{PROV-O.}
The W3C PROV family~\cite{w3c2013provo} supplies
\texttt{prov:Entity}, \texttt{prov:Activity}, \texttt{prov:Agent}, and
the relations we lean on (\texttt{prov:wasGeneratedBy},
\texttt{prov:wasAttributedTo}, \texttt{prov:wasDerivedFrom},
\texttt{prov:hadPlan}). We add a small project namespace
(\texttt{fair2r:}) with subclasses for the new artefact classes:
\texttt{fair2r:Claim}, \texttt{fair2r:Source}, \texttt{fair2r:Prompt}
(a \texttt{prov:Plan}), \texttt{fair2r:AIAgent} and
\texttt{fair2r:HumanResearcher}, \texttt{fair2r:AuthoringPass} and
\texttt{fair2r:AuditPass}, and a \texttt{fair2r:VerificationState}
class whose individuals are the rungs of the ladder defined in
\S\ref{sec:pattern}. PROV-O is necessary but not sufficient: it
records that an activity happened, not how confidently a claim was
verified. That is the ladder's job.

\paragraph{Retrieval depth and certainty of evidence.}
PRISMA 2020~\cite{page2021prisma} provides a retrieval-depth axis
(Identification, Screening, Included). GRADE~\cite{guyatt2008grade}
provides an invocation-strength axis (very low, low, moderate, high).
Neither was designed for LLM-assisted writing. The verification
ladder of \S\ref{sec:pattern} is best read as a crosswalk between
them: PRISMA gives the floor a rung permits; GRADE rates the
underlying evidence, not the retrieval procedure.

\paragraph{Disclosure norms.}
The \textsc{ICMJE} Recommendations and most major venues now require
authors to disclose the use of generative AI and to take
responsibility for the output~\cite{icmje2023}. All converge on one
position: AI is not an author. The legal background is consistent ---
in Germany (\S2 \emph{UrhG}~\cite{urhg2}\todo[inline]{verify}), in the
United States (USCO 2023~\cite{usco2023ai}\todo[inline]{verify}), and
through the major medical journals --- AI cannot hold copyright,
cannot consent to publication, and cannot be held accountable. The
disclosure landscape covers \emph{whether} to disclose; it does not
yet cover \emph{what to preserve}. F(AI)\textsuperscript{2}R is the
proposal that there is a usable answer.
